\chapter{Fourier Transform Properties}

\section{Symmetry}

\subsection{Even and Odd Functions}

To appreciate the Symmetry Property of the Fourier transform, we must understand the decomposition of functions into even and odd components. Any function $f(t)$ has a unique decomposition into even and odd functions $f_e$ and $f_o$, where

\begin{equation}
	f_e(t) = \frac{f(t) + f(-t)}{2}
\end{equation}

and

\begin{equation}
	f_o(t) = \frac{f(t) - f(-t)}{2}
\end{equation}

\subsection{The Symmetry Property}

We can decompose real and imaginary parts of a complex function into even and odd sections to give four components (denoted below as RE, RO, IE, and IO). The Symmetry Property of the Fourier transform states that there is a one-to-one mapping between the four components of a complex time function and the four components of its complex frequency transform. That mapping is shown below.

\vspace{0.5cm}

\begin{tikzpicture}[yscale=1.5, xscale=1.8]
	% --- Row 1: Time Domain ---
	\node at (-1.5, 1) {\textbf{Time domain}};
	\node (f)   at (0, 1) {$f$};
	\node (eq1) at (0.5, 1) {$=$};
	\node (fRE) at (1, 1) {$f_{\text{RE}}$};
	\node (p1)  at (1.4, 1) {$+$};
	\node (fRO) at (1.8, 1) {$f_{\text{RO}}$};
	\node (p2)  at (2.2, 1) {$+$};
	\node (ifIE) at (2.6, 1) {$i\,f_{\text{IE}}$};
	\node (p3)  at (3, 1) {$+$};
	\node (ifIO) at (3.4, 1) {$i\,f_{\text{IO}}$};
	
	% --- Row 2: Arrows ---
	% Custom double arrow with F symbol
	\newcommand{\farrow}[1]{
		\draw[<->, >=Latex] (#1, 0.8) -- (#1, 0.2) 
		node[midway, right, xshift=2pt] {$\mathcal{F}$};
	}
	
	\farrow{0} \farrow{1} \farrow{1.8} \farrow{2.6} \farrow{3.4}
	
	% --- Row 3: Frequency Domain ---
	\node at (-1.5, 0) {\textbf{Frequency domain}};
	\node (hf)   at (0, 0) {$F$};
	\node (eq2)  at (0.5, 0) {$=$};
	\node (hfRE) at (1, 0) {$F_{\text{RE}}$};
	\node (p4)   at (1.4, 0) {$+$};
	\node (ifIOh) at (1.8, 0) {$i\,F_{\text{IO}}$};
	\node (p5)   at (2.2, 0) {$+$};
	\node (ifIEh) at (2.6, 0) {$i\,F_{\text{IE}}$};
	\node (p6)   at (3, 0) {$+$};
	\node (hfRO) at (3.4, 0) {$F_{\text{RO}}$};
	
\end{tikzpicture}

\vspace{0.5cm}

From the Symmetry Property, numerous relationships are apparent. Just a few are summarized below.

\subsection{Real Functions}

If $f(t)$ is real, then its Fourier transform will have an even real component and an odd imaginary component. This type of function is called Hermitian or conjugate symmetric and satisfies the condition below.

\begin{equation}
	F(-\omega) = F^*(\omega)
\end{equation}

Conversely, if the Fourier transform of $f(t)$ is conjugate symmetric, then $f(t)$ is real.

\subsection{Imaginary Functions}

If $f(t)$ is imaginary, then its Fourier transform will have an odd real component and an even imaginary component. This type of function is called anti-Hermitian or conjugate anti-symmetric and satisfies the condition below.

\begin{equation}
	F(-\omega) = -F^*(\omega)
\end{equation}

Conversely, if the Fourier transform of $f(t)$ is conjugate anti-symmetric, then $f(t)$ is imaginary.

\subsection{Conjugate Symmetric Functions}
If $f(t)$ is conjugate symmetric (i.e. $f(t) = f_{RE}(t) + f_{IO}(t)$), then its Fourier transform will be real. Conversely, if the Fourier transform of $f(t)$ is real then $f(t)$ is conjugate symmetric.

\subsection{Conjugate Anti-Symmetric Functions}
If $f(t)$ is conjugate anti-symmetric (i.e. $f(t) = f_{RO}(t) + f_{IE}(t)$), then its Fourier transform will be imaginary. Conversely, if the Fourier transform of $f(t)$ is imaginary then $f(t)$ is conjugate anti-symmetric.

\section{Hilbert Transform Relations}

\subsection{Introduction}
In system theory, Hilbert Transform Relations are a set of unique relationships between the real and imaginary parts of a complex function resulting from various constraints.

\subsection{Causal Real Functions}
\textbf{Derivation}: Let $f$ be a causal, stable, real function with Fourier transform $F(\omega)$. Let $f_e$ and $f_o$ be the even and odd components of $f$ respectively. We seek to express $f$ in terms of $f_e$:

\begin{gather*}
	f_e(t) = \frac{f(t) + f(-t)}{2} \\
	2f_e(t) = f(t) + f(-t) \\
	2f_e(t)u(t) = u(t)(f(t) + f(-t))
\end{gather*}

Because $f$ is causal, $f(t)$ is zero for all $t < 0$ and $f(-t)$ is zero for all $t > 0$. It follows that $f(t)u(t) = f(t)$ and $f(-t)u(t) = f(0)\delta(t)$. Then

\begin{equation*}
	2f_e(t)u(t) = f(t) + \delta(t)f(0)
\end{equation*}

Since $f(0) = f_e(0)$, we can express $f$ entirely in terms of $f_e$.

\begin{equation} \label{eq:causal_function_in_terms_of_even_component}
	f(t) = 2u(t)f_e(t) - \delta(t)f_e(0)
\end{equation}

Recall that by the Symmetry Property, the Fourier transform of $f_e$ is real and that of $f_o$ is imaginary. It follows that $F$ can be determined from its real component $F_R$.

\textbf{Summary}: The Fourier transform of a causal, stable, real function can be uniquely determined from its real component by the following process:

\begin{enumerate}
	\item Use the inverse Fourier transform to find $f_e$ from $F_R$.
	\item Recover $f$ from $f_e$ using \eqref{eq:causal_function_in_terms_of_even_component}.
	\item Use the Fourier transform to find $F$ from $f$.
\end{enumerate}

It follows, of course, that $F_I$ can be determined from $F_R$.



\subsection{Finite Causal Real Discrete-Time Functions}
Let $f$ be a causal, finite, real, discrete-time function with DFT $F[k]$. 

\textbf{Summary}: The discrete Fourier transform of a causal, finite, real function can be determined uniquely from its real component by the following process:

\begin{enumerate}
	\item Compute the inverse DFT of $F_R[k]$ to obtain the sequence
	\begin{equation*}
		f_{ep}[n] = \frac{f[n] + f[((-n))_N]}{2}, 0 \leq n \leq N - 1.
	\end{equation*}
	
	\item Compute the periodic odd part of $f[n]$ as
	\begin{equation*}
		f_{op}[n] = \begin{cases}
			f_{ep}[n], & 0 < n < \frac{N}{2}, \\
			-f_{ep}[n], & \frac{N}{2} < n \leq N - 1, \\
			0, & otherwise.
		\end{cases}
	\end{equation*}
	
	\item Compute the DFT of $f_{op}$ to obtain $jF_I[k]$.
\end{enumerate}

\subsection{Complex Sequences}
Let $x$ be a complex sequence with a Fourier transform $X$ such that

\begin{equation*}
	X(e^{j\omega}) = 0, -\pi \leq \omega < 0.
\end{equation*}

Let $x_r$ and $x_i$ be the real and imaginary components of $x$ respectively such that $x[n] = x_r[n] + ix_i[n]$. Let $X_r$ and $X_i$ be the Fourier transforms of $x_r$ and $x_i$ respectively such that $X(e^{j\omega}) = X_r(e^{j\omega}) + X_i(e^{j\omega})$ (note that $X_r$ and $X_i$ are \textit{not} the real and imaginary components of $X$). Then

\begin{equation}
	X(e^{j\omega}) = \begin{cases}
		2X_r(e^{j\omega}), 0 < \omega < \pi \\
		0, -\pi \leq \omega < 0,
	\end{cases}
\end{equation}

and

\begin{equation}
	X(e^{j\omega}) = \begin{cases}
		2jX_i(e^{j\omega}), 0 < \omega < \pi \\
		0, -\pi \leq \omega < 0.
	\end{cases}
\end{equation}

Note that $X$ cannot be recovered from either $X_r$ or $X_i$ alone at $\omega = 0$.

Additionally,

\begin{equation}
	X_i(e^{j\omega}) = H(e^{j\omega})X_r(e^{j\omega}),
\end{equation}

and

\begin{equation}
	X_r(e^{j\omega}) = -H(e^{j\omega})X_i(e^{j\omega}),
\end{equation}

where

\begin{equation}
	H(e^{j\omega}) = \begin{cases}
		-j, 0 < \omega < \pi, \\
		j, \pi < \omega < 0.
	\end{cases}
\end{equation}

A system with frequency response $H$ is called an ideal 90-degree phase shifter, or a Hilbert Transformer.

